\documentclass{beamer}

\usepackage{cuhkbeamer}

\begin{document}

% ====================== Slide 1: TITLE SLIDE ======================

\begin{frame}
  \title{Cost Structures and Revenue Management}
  \author{Chiu Yu KO}
  \institute{Low Altitude Economics}
  \date{Week 3}
  \titlepage
\end{frame}

% ====================== Slide 2: FORMAT & OBJECTIVES ======================
\begin{frame}[t]
  \frametitle{Learning Objectives}
  By the end of this class, you should be able to:
  \begin{itemize}
    \item \textbf{Distinguish} CAPEX from OPEX in an LAE business model.
    \item \textbf{Calculate} simple unit-cost measures, including CASK and cost per mission-kilometre.
    \item \textbf{Estimate} breakeven volume and explain how price, load factor, and utilization affect viability.
    \item \textbf{Compare} pay-per-use, subscription, availability-based, and complementary revenue models.
    \item \textbf{Evaluate} how weather downtime and cost differences across Hong Kong and the GBA affect business strategy.
  \end{itemize}
\end{frame}

% ====================== Slide 3: ICEBREAKER ======================
\begin{frame}[t]
  \frametitle{Midnight in Kowloon: The Numbers Behind the Service}
  \textbf{Imagine two managers reviewing an LAE business case.}
  \begin{itemize}
    \item \textbf{Manager A} asks how charging time, maintenance, and turnaround affect the number of daily missions.
    \item \textbf{Manager B} asks whether subscriptions or availability-based contracts could make revenue more predictable.
  \end{itemize}
  \vspace{0.3cm}
  A technically feasible service is not necessarily commercially viable.
  \begin{itemize}
    \item Managers must connect asset cost, daily operations, capacity, demand, pricing, and risk.
    \item \textbf{Quick poll (2 min):} What is likely to be the most important cost driver for an LAE operator in Hong Kong?
  \end{itemize}
\end{frame}

% ====================== Session 1 Title (after Icebreaker) ======================

\begin{frame}[c]
  \begin{center}
    \huge\color{cuhkpurple}\textbf{Session 1}\\
    \vspace{0.5cm}
    \Large CAPEX vs OPEX, Unit Economics \& Cost Drivers
  \end{center}
\end{frame}

% ====================== Slide 4: CONCEPT PRIMER 1 ======================
\begin{frame}[t]
  \frametitle{Concept Primer: CAPEX and OPEX}
  \textbf{CAPEX: Investment in assets that provide benefits over several years}
  \begin{itemize}
    \item Examples: aircraft, charging equipment, vertiport construction, and major software systems.
  \end{itemize}
  \textbf{OPEX: Recurring expenditure required to operate the service}
  \begin{itemize}
    \item Examples: energy, maintenance, staff, insurance, site fees, communications, and routine compliance activities.
  \end{itemize}
  \vspace{0.2cm}
  \textbf{Important distinction:} Certification and approval expenditure is not automatically CAPEX. Its accounting treatment depends on the nature of the expenditure and applicable accounting rules.
\end{frame}

% ====================== Slide 5: THE HK PREMIUM ======================
\begin{frame}[t]
  \frametitle{The ``Hong Kong Premium'': Potential Cost Drivers}
  Why might an LAE service cost more to operate in Hong Kong?
  \begin{itemize}
    \item \textbf{Land and site adaptation:} Scarce sites, structural reinforcement, charging capacity, passenger access, and safety facilities may raise infrastructure cost.
    \item \textbf{Complex operating environment:} Dense buildings, terrain, communications, weather, and interactions with existing airspace users may increase operating requirements.
    \item \textbf{Labour and professional services:} Skilled staff, engineering, insurance, legal, and compliance support may be costly.
    \item \textbf{Reliability requirements:} Redundancy, maintenance capability, backup systems, and contingency planning can add cost.
  \end{itemize}
  \textit{These are cost hypotheses to test. The size of any Hong Kong premium will depend on the application, site, and operating model.}
\end{frame}

% ====================== NEW SLIDE A: APPLIED TOOL - CAPEX + OPEX ======================
\begin{frame}[t]
  \frametitle{Applied Tool: Estimating Daily Economic Cost}
  \begin{block}{A Simple Managerial Approximation}
  \[
    \text{Daily Economic Cost}
    =
    \frac{\text{Asset Cost}-\text{Expected Residual Value}}{\text{Useful Operating Days}}
    +\text{Daily OPEX}
  \]
  \end{block}
  \begin{itemize}
    \item This spreads the net asset cost across its expected operating life.
    \item Daily OPEX includes both fixed operating costs and costs that vary with activity.
    \item Financing cost, maintenance cycles, asset downtime, and residual value may require separate treatment in a full model.
  \end{itemize}
  \textit{This is an economic-cost approximation, not a complete accounting depreciation or cash-flow model.}
\end{frame}

% ====================== NEW SLIDE B: BACK-OF-THE-ENVELOPE (CAPEX + OPEX) ======================
\begin{frame}[t,shrink=12]
  \frametitle{Back-of-the-Envelope: How Asset and Site Cost Affect Unit Cost}
  \textbf{Assumptions:} Five-year life, no residual value, daily OPEX of HK\$8,000, and 50 missions per operating day.
  \begin{columns}[T]
    \begin{column}{0.48\textwidth}
      \begin{block}{Case A: HK\$10 Million Investment}
      \begin{itemize}
        \item Daily asset-cost allocation: HK\$10M / 1,825 $\approx$ HK\$5,479
        \item Daily OPEX: HK\$8,000
        \item Total daily economic cost: HK\$13,479
        \item \textbf{Cost per mission: about HK\$270}
      \end{itemize}
      \end{block}
    \end{column}
    \begin{column}{0.48\textwidth}
      \begin{block}{Case B: HK\$6 Million Investment}
      \begin{itemize}
        \item Daily asset-cost allocation: HK\$6M / 1,825 $\approx$ HK\$3,288
        \item Daily OPEX: HK\$8,000
        \item Total daily economic cost: HK\$11,288
        \item \textbf{Cost per mission: about HK\$226}
      \end{itemize}
      \end{block}
    \end{column}
  \end{columns}
  \textit{Takeaway: Higher investment cost raises unit cost, but utilization and operating cost remain equally important.}
  \vfill
  {\tiny Illustrative assumptions for classroom analysis; not current market quotations.}
\end{frame}


% ====================== Slide 6: CONCEPT PRIMER 2 ======================
\begin{frame}[t]
  \frametitle{Concept Primer: ``Unit Economics''}
  Unit economics examines the revenue and cost associated with one meaningful unit of activity.
  \begin{itemize}
    \item Passenger service: per flight, available seat-kilometre, or paying passenger-trip.
    \item Cargo or public-service operation: per mission, mission-kilometre, or kilogram delivered.
  \end{itemize}
  \vspace{0.25cm}
  \begin{itemize}
    \item \textbf{CASK}: Operating cost divided by available seat-kilometres. It measures the cost of capacity offered, whether seats are occupied or not.
    \item \textbf{Cost per mission-kilometre}: A simple mission-based measure for comparing cargo, inspection, and public-service operations.
  \end{itemize}
  \textit{Different metrics answer different managerial questions and should be interpreted together with demand and service quality.}
\end{frame}

% ====================== Slide 7: STANDALONE FORMULA ======================
\begin{frame}[t]
  \frametitle{Applied Tool: Calculating CASK}
  \begin{block}{Cost per Available Seat-Kilometre}
  \[
    CASK=\frac{\text{Operating Cost for the Period}}{\text{Available Seats}\times\text{Kilometres Flown}}
  \]
  \end{block}
  \begin{itemize}
    \item Available seat-kilometres measure the passenger capacity offered, not the number of seats sold.
    \item More productive aircraft use can spread some fixed operating costs over more seat-kilometres.
    \item A low CASK does not guarantee profitability. The operator must also fill enough seats and earn sufficient revenue.
    \item Turnaround, charging, maintenance, weather, route length, and available capacity all affect the measure.
  \end{itemize}
\end{frame}

% ====================== Slide 8: BACK-OF-ENVELOPE (CASK) ======================
\begin{frame}[t,shrink=10]
  \frametitle{Back-of-the-Envelope: How Utilization Affects CASK}
  \textbf{Assumptions:} Four available seats and total operating cost of HK\$12,000 per day.
  \begin{columns}[T]
    \begin{column}{0.48\textwidth}
      \begin{block}{10 Trips of 10 km}
      \begin{itemize}
        \item Available seat-km: $4\times10\times10=400$
        \item $CASK=12{,}000/400$
        \item \textbf{CASK = HK\$30 per seat-km}
        \item Allocated cost of one available seat on a 10 km trip: HK\$300
      \end{itemize}
      \end{block}
    \end{column}
    \begin{column}{0.48\textwidth}
      \begin{block}{50 Trips of 10 km}
      \begin{itemize}
        \item Available seat-km: $4\times50\times10=2{,}000$
        \item $CASK=12{,}000/2{,}000$
        \item \textbf{CASK = HK\$6 per seat-km}
        \item Allocated cost of one available seat on a 10 km trip: HK\$60
      \end{itemize}
      \end{block}
    \end{column}
  \end{columns}
  \textit{Takeaway: Higher utilization lowers CASK in this simplified example, but does not establish market demand or profitability.}
  \vfill
  {\tiny Illustrative assumptions for classroom analysis.}
\end{frame}


% ====================== NEW SLIDE: APPLIED TOOL - CPMK ======================
\begin{frame}[t]
  \frametitle{Applied Tool: Cost per Mission-Kilometre}
  \begin{block}{A Mission-Based Cost Measure}
  \[
    \text{Cost per Mission-Km}
    =
    \frac{\text{Operating Cost for the Period}}{\text{Total Mission-Kilometres}}
  \]
  \end{block}
  \begin{itemize}
    \item Useful for comparing cargo, inspection, emergency, or public-service operations.
    \item Lower cost can come from reducing operating cost or increasing productive mission volume.
    \item Longer missions mechanically create more mission-kilometres, so comparisons should consider route length, payload, service quality, and mission complexity.
  \end{itemize}
\end{frame}

% ====================== NEW SLIDE: BACK-OF-THE-ENVELOPE - CPMK ======================
\begin{frame}[t,shrink=10]
  \frametitle{Back-of-the-Envelope: Mission Volume and Unit Cost}
  \textbf{Assumptions:} Daily operating cost = HK\$12,000; average mission distance = 8 km.
  \begin{columns}[T]
    \begin{column}{0.48\textwidth}
      \begin{block}{20 Missions per Day}
      \begin{itemize}
        \item Mission-km: $20\times8=160$
        \item Cost per mission-km: HK\$75
        \item \textbf{Allocated cost per mission: HK\$600}
      \end{itemize}
      \end{block}
    \end{column}
    \begin{column}{0.48\textwidth}
      \begin{block}{100 Missions per Day}
      \begin{itemize}
        \item Mission-km: $100\times8=800$
        \item Cost per mission-km: HK\$15
        \item \textbf{Allocated cost per mission: HK\$120}
      \end{itemize}
      \end{block}
    \end{column}
  \end{columns}
  \textit{Takeaway: Greater productive volume can reduce allocated unit cost, but capacity, demand, turnaround, and service requirements constrain volume.}
  \vfill
  {\tiny Illustrative assumptions for classroom analysis.}
\end{frame}

% ====================== Slide 9: SESSION 2 TITLE ======================
\begin{frame}[c]
  \begin{center}
    \huge\color{cuhkpurple}\textbf{Session 2}\\
    \vspace{0.5cm}
    \Large Revenue Streams, Yield Management \& Breakeven
  \end{center}
\end{frame}

% ====================== Slide 10: REVENUE DIVERSIFICATION ======================
\begin{frame}[t]
  \frametitle{Moving Beyond a Single Transaction}
  LAE businesses may combine several revenue models:
  \begin{enumerate}
    \item \textbf{Pay per use:} The customer pays for each passenger trip, delivery, inspection, or mission.
    \item \textbf{Subscription or membership:} The customer pays periodically for access, priority, discounts, or a defined service allowance.
    \item \textbf{Availability-based contract:} A public agency or business pays for readiness and agreed service levels, with possible additional usage payments.
    \item \textbf{Complementary services:} Data, analytics, maintenance, platform, or integration services may create additional revenue, subject to consent, data rights, privacy, and customer demand.
  \end{enumerate}
\end{frame}

% ====================== NEW SLIDE: DESIGNING REVENUE STREAMS ======================
% ====================== NEW SLIDE: DESIGNING REVENUE STREAMS ======================
\begin{frame}[t]
  \frametitle{Choosing a Revenue Model}
  \begin{center}
  \small
  \begin{tabular}{l|ccc}
    \textbf{Model} & Revenue predictability & Usage risk & Possible application\\
    \hline
    Pay per use & Lower & Operator & Passenger or delivery\\
    Subscription & Medium & Shared & Frequent users\\
    Availability contract & Higher & Customer / shared & Public or critical service\\
    Complementary service & Varies & Varies & Data or platform support\\
  \end{tabular}
  \end{center}
  \vspace{0.3cm}
  \begin{itemize}
    \item Match the revenue model to the value promised and the party best able to manage the risk.
    \item A subscription is not guaranteed revenue unless customers accept the contract terms and continue paying.
    \item Data does not generate revenue automatically. It requires useful information, data rights, customer demand, security, and privacy compliance.
    \item Combining models may diversify revenue but also increases contractual and operational complexity.
  \end{itemize}
\end{frame}
% ====================== Slide 11: YIELD MANAGEMENT ======================
\begin{frame}[t]
  \frametitle{Concept Primer: Revenue Management}
  Revenue management allocates limited capacity and adjusts prices or conditions to improve total expected revenue.
  \begin{itemize}
    \item The objective is not simply to eliminate every empty seat.
    \item Selling too cheaply may fill capacity while reducing total revenue.
    \item Holding capacity for later high-value demand can sometimes be worthwhile.
    \item Managers consider price, booking time, customer segment, expected demand, cancellation risk, and capacity.
  \end{itemize}
  \textbf{Managerial objective:} Sell the right capacity to the right customer at the right time and price.
\end{frame}


% ====================== Slide 12: BREAKEVEN EQUATION ======================
\begin{frame}[t]
  \frametitle{Applied Tool: The Breakeven Equation}
  \begin{block}{One Service Unit}
  \[
    \text{Breakeven Units}
    =
    \frac{\text{Fixed Cost}}{\text{Revenue per Unit}-\text{Variable Cost per Unit}}
  \]
  \end{block}
  \begin{itemize}
    \item The denominator is the \textbf{contribution margin per unit}.
    \item The unit must be consistent: one flight, one passenger, one delivery, or one mission.
    \item For passenger operations, revenue per flight depends on both average fare and paying passengers per flight.
  \end{itemize}
\end{frame}

% ====================== Slide 13: BREAKEVEN CALCULATION ======================
\begin{frame}[t,shrink=10]
  \frametitle{Back-of-the-Envelope: Price and Breakeven Volume}
  \textbf{Assumptions:} Fixed cost = HK\$20,000/day; variable cost = HK\$100 per mission.
  \begin{columns}[T]
    \begin{column}{0.48\textwidth}
      \begin{block}{Revenue of HK\$500 per Mission}
      \begin{itemize}
        \item Contribution margin: HK\$400
        \item Breakeven missions: $20{,}000/400$
        \item \textbf{50 missions per day}
      \end{itemize}
      \end{block}
    \end{column}
    \begin{column}{0.48\textwidth}
      \begin{block}{Revenue of HK\$200 per Mission}
      \begin{itemize}
        \item Contribution margin: HK\$100
        \item Breakeven missions: $20{,}000/100$
        \item \textbf{200 missions per day}
      \end{itemize}
      \end{block}
    \end{column}
  \end{columns}
  \textit{Takeaway: Discounting can sharply increase required volume. The manager must test whether the target fits available capacity and expected demand.}
  \vfill
  {\tiny Illustrative assumptions for classroom analysis.}
\end{frame}

% ====================== Slide 14: UTILIZATION RATE ======================
\begin{frame}[t]
  \frametitle{Applied Tool: Required Daily Activity}
  \begin{block}{A Breakeven Activity Target}
  \[
    \text{Required Missions per Day}
    =
    \frac{\text{Fixed Daily Cost}}{\text{Contribution Margin per Mission}}
  \]
  \end{block}
  \begin{itemize}
    \item Compare the required activity with physical capacity, operating hours, charging, maintenance, weather, and expected demand.
    \item This is not the usual definition of a utilization rate. A utilization rate is normally actual productive use divided by available capacity or time.
    \item Revenue management can improve contribution, but cannot create unlimited capacity or demand.
  \end{itemize}
\end{frame}

% ====================== Slide 15: YIELD MANAGEMENT AS THE TOOL ======================
\begin{frame}[t]
  \frametitle{Revenue Management and Revenue per Flight}
  \begin{block}{A Simple Passenger-Revenue Relationship}
  \[
    \text{Revenue per Flight}
    =
    \text{Seats}\times\text{Load Factor}\times\text{Average Fare}
  \]
  \end{block}
  \begin{itemize}
    \item Load factor is the proportion of available seats occupied by paying passengers.
    \item Average fare should reflect discounts, customer mix, and booking conditions.
    \item The objective is to improve total contribution, not merely to maximize load factor.
  \end{itemize}
\end{frame}


% ====================== Slide 15: BACK-OF-THE-ENVELOPE (YIELD MANAGEMENT) ======================
\begin{frame}[t,shrink=15]
  \frametitle{Back-of-the-Envelope: Load Factor and Breakeven}
  \textbf{Assumptions:} Four seats, average fare HK\$500, fixed cost HK\$20,000/day, and variable cost HK\$100/flight.
  \begin{columns}[T]
    \begin{column}{0.48\textwidth}
      \begin{block}{40\% Load Factor}
      \begin{itemize}
        \item Expected paying passengers: $4\times40\%=1.6$
        \item Revenue per flight: $1.6\times500=\text{HK\$800}$
        \item Contribution: HK\$700
        \item Breakeven: $20{,}000/700\approx\mathbf{29}$ flights/day
      \end{itemize}
      \end{block}
    \end{column}
    \begin{column}{0.48\textwidth}
      \begin{block}{90\% Load Factor}
      \begin{itemize}
        \item Expected paying passengers: $4\times90\%=3.6$
        \item Revenue per flight: $3.6\times500=\text{HK\$1,800}$
        \item Contribution: HK\$1,700
        \item Breakeven: $20{,}000/1{,}700\approx\mathbf{12}$ flights/day
      \end{itemize}
      \end{block}
    \end{column}
  \end{columns}
  \textit{Takeaway: Higher load factor can reduce the required number of flights, provided fares and costs remain unchanged.}
  \vfill
  {\tiny Illustrative assumptions; fractional passengers represent an average across many flights.}
\end{frame}

% ====================== Slide 15: SESSION 3 TITLE ======================
\begin{frame}[c]
  \begin{center}
    \huge\color{cuhkpurple}\textbf{Session 3}\\
    \vspace{0.5cm}
    \Large Cost Sensitivities, Risk Hedging \& GBA Scale
  \end{center}
\end{frame}

% ====================== Slide 16: CONCEPT PRIMER 4 ======================
\begin{frame}[t]
  \frametitle{Concept Primer: Weather Downtime}
  Severe weather can affect an LAE business through several channels:
  \begin{itemize}
    \item lost trips and revenue;
    \item continued aircraft, site, staff, insurance, and financing costs;
    \item customer disruption, refunds, and service-level penalties;
    \item repositioning, inspection, maintenance, and recovery costs; and
    \item lower customer willingness to rely on the service.
  \end{itemize}
  \textit{A warning signal does not mechanically determine every aviation operation. Actual restrictions depend on aircraft capability, location, operational approval, and prevailing conditions.}
\end{frame}

% ====================== Slide 17: CONCEPT PRIMER 5 ======================
\begin{frame}[t]
  \frametitle{Concept Primer: Managing Weather Risk}
  Managers can combine operational, contractual, and financial responses:
  \begin{itemize}
    \item \textbf{Operational resilience:} Forecasting, spare capacity, alternative routes or modes, and recovery planning.
    \item \textbf{Contract design:} Availability payments, minimum commitments, cancellation terms, service credits, and force-majeure provisions.
    \item \textbf{Insurance or parametric cover:} A payment may be triggered by an agreed event or weather measure, subject to availability, exclusions, basis risk, and premium cost.
    \item \textbf{Liquidity planning:} Cash reserves or committed financing to absorb temporary revenue shortfalls.
  \end{itemize}
  \textit{Subscriptions and reservation fees transfer risk only if customers accept the terms and receive a credible value proposition.}
\end{frame}

% ====================== Slide 18: STANDALONE FORMULA ======================
\begin{frame}[t]
  \frametitle{Applied Tool: Profit Under a Mixed-Revenue Model}
  \begin{block}{A Simple Monthly Model}
  \[
  \begin{split}
    \text{Profit}={}&\text{Committed Revenue}+\text{Usage Revenue}+\text{Insurance Payout}\\
    &-\text{Fixed Cost}-\text{Variable Cost}-\text{Premiums and Service Credits}
  \end{split}
  \]
  \end{block}
  \begin{itemize}
    \item Committed revenue may reduce exposure to flight volume, but often comes with availability or service obligations.
    \item Downtime may save some variable cost while still creating refunds, credits, recovery costs, or reputational damage.
    \item The objective is not to profit from disruption. It is to keep cash flow and service obligations manageable.
  \end{itemize}
\end{frame}

% ====================== Slide 19: BACK-OF-ENVELOPE (HEDGING) ======================
\begin{frame}[t,shrink=15]
  \frametitle{Back-of-the-Envelope: A Month with Five Grounded Days}
  \textbf{Assumptions:} 30 days; 25 operating days; 50 missions/day; HK\$500 usage revenue; fixed cost HK\$600,000/month; variable cost HK\$100/mission.
  \begin{columns}[T]
    \begin{column}{0.48\textwidth}
      \begin{block}{Pay per Mission}
      \begin{itemize}
        \item Missions: 1,250
        \item Revenue: HK\$625,000
        \item Variable cost: HK\$125,000
        \item Fixed cost: HK\$600,000
        \item \textbf{Operating loss: HK\$100,000}
      \end{itemize}
      \end{block}
    \end{column}
    \begin{column}{0.48\textwidth}
      \begin{block}{Committed Monthly Revenue}
      \begin{itemize}
        \item Contract revenue: HK\$750,000
        \item Variable cost: HK\$125,000
        \item Fixed cost: HK\$600,000
        \item \textbf{Operating profit before service credits: HK\$25,000}
      \end{itemize}
      \end{block}
    \end{column}
  \end{columns}
  \textit{Takeaway: Committed revenue may stabilize cash flow, but the contract must specify availability, refunds, service credits, and weather-related obligations.}
  \vfill
  {\tiny Illustrative assumptions; excludes taxes, financing, insurance premiums, recovery costs, and possible service credits.}
\end{frame}

% ====================== Slide 20: HK VS CHINA VARIATIONS ======================
\begin{frame}[t]
  \frametitle{Hong Kong and the GBA: Different Cost and Revenue Possibilities}
  \begin{center}
  \small
  \begin{tabular}{l|cc}
    \textbf{Dimension} & \textbf{Hong Kong possibility} & \textbf{Wider GBA possibility}\\
    \hline
    Site environment & Dense and constrained & More varied\\
    Customer base & Premium and specialized & Larger and more diverse\\
    Cost pressure & Site and professional costs & Varies by city and use case\\
    Scale opportunity & Complex urban services & Manufacturing and network scale\\
    Potential role & Finance, logistics, services & Technology, production, operations\\
  \end{tabular}
  \end{center}
  \vspace{0.35cm}
  \textit{These are strategic hypotheses rather than universal facts. Smart operators should compare specific sites, routes, customers, regulations, and partners.}
\end{frame}

% ====================== Slide 19: EXERCISE INSTRUCTIONS ======================
\begin{frame}[t]
  \frametitle{In-Class Exercise: Spreadsheet Workshop}
  
  \textbf{Activity: Unit Economics \& Sensitivity Testing}
  \begin{itemize}
    \item \textbf{Format:} Groups of 4–5 (Laptops required)
    \item \textbf{Scenario:} You are managing a new HK Airport Shuttle fleet.
    \item \textbf{Task:} Using the provided Excel template:
    \begin{enumerate}
        \item Input the Base Case assumptions to find your Breakeven.
        \item \textbf{Shock 1:} Increase Vertiport Rent by 20\%.
        \item \textbf{Shock 2:} Simulate 3 days of Typhoon downtime.
        \item How much do you have to raise ticket prices to survive?
    \end{enumerate}
    \item \textbf{Timing:} 15 minutes hands-on + 5 minute share-out.
  \end{itemize}
\end{frame}

% ====================== Slide 20: EXERCISE TEMPLATE ======================
\begin{frame}[t]
  \frametitle{Exercise Template (Visual)}
  
  \begin{center}
      \textit{[Distribute digital Excel Spreadsheet to student groups]}
      
      \vspace{0.5cm}
      
      \textbf{The Dashboard View:} \\
      \vspace{0.2cm}
      Input: Fleet Size | Daily Rent | Electricity Cost | Ticket Price \\
      Output: Daily Profit | Breakeven Flights | CASK \\
      \vspace{0.5cm}
      \textit{Watch how the "Daily Profit" box turns red the second you lower the utilisation rate!}
  \end{center}
\end{frame}

% ====================== Slide 21: KEY TAKEAWAYS ======================
\begin{frame}[t]
  \frametitle{Key Takeaways}
  \begin{itemize}
    \item Separate CAPEX, recurring fixed OPEX, and activity-related variable cost before calculating unit economics.
    \item CASK measures the cost of passenger capacity offered; mission-based measures may be more useful for cargo and public-service applications.
    \item Breakeven depends on a consistent unit, contribution margin, feasible capacity, and realistic demand.
    \item Revenue management aims to improve total contribution, not simply maximize load factor.
    \item Subscriptions, availability contracts, data services, insurance, and liquidity can reshape risk, but none provides a free or guaranteed hedge.
    \item Hong Kong and other GBA cities may offer complementary cost, customer, technology, and service advantages.
  \end{itemize}
\end{frame}

% ====================== Slide 22: ASSIGNMENT ======================
\begin{frame}[t]
  \frametitle{Week 3 Assignment Briefing}
  \textbf{Cost and Revenue Model with Memo (maximum 1,200 words plus spreadsheet)}
  \begin{itemize}
    \item \textbf{Due:} Beginning of Week 4 (group assignment)
    \item \textbf{Task:} Extend the Week 2 corridor or application into a simple operating model.
    \item \textbf{Requirements:}
    \begin{enumerate}
      \item Separate investment cost, fixed OPEX, and variable cost.
      \item Calculate one appropriate unit-cost metric and breakeven volume.
      \item Model revenue per mission using price, capacity, and expected utilization or load factor.
      \item Test three shocks: lower price, weather downtime, and higher compliance or infrastructure cost.
      \item Recommend one operational, contractual, or financial risk response and discuss its limitation.
    \end{enumerate}
    \item \textbf{Grading focus:} Accurate and internally consistent model (40\%), transparent assumptions (30\%), actionable recommendation (20\%), professionalism (10\%).
  \end{itemize}
\end{frame}

% ====================== Slide 23: PREVIEW WEEK 4 ======================
\begin{frame}[t]
  \frametitle{Preview of Week 4}
  
  \begin{center}
      \Large \textbf{Demand Modeling in Urban Air Mobility}
  \end{center}
  
  \vspace{1cm}
  \textbf{Preview question for next week:} \\
  \textit{Now that we know exactly how much we HAVE to charge to survive... how do we mathematically prove that customers are actually willing to pay it?}
\end{frame}

% ====================== Slide 24: THANK YOU ======================
\begin{frame}[t]
  \frametitle{Thank You + Q\&A}
  
  \textbf{Pre-Class Readings Reminder (for Week 4):}
  \begin{itemize}
      \item \textit{McKinsey: Perspectives on AAM (Adoption patterns)}
      \item \textit{NASA UAM Market Study (Trip generation \& Mode choice)}
  \end{itemize}
  
  \vspace{1cm}
  \begin{center}
      \textbf{Questions?} \\
      \small (Final 15 minutes reserved for extended Q\&A)
  \end{center}
\end{frame}

\end{document}